\chapter{Validity: Are the Two Systems Comparable?}
\label{ch:validity}

% SOURCE OF TRUTH: thesis-pack/08-validity/PARITY.md (authoritative on every
% parity claim and every deviation). Bound to workflow_baseline/test_parity.py.
% Supporting: ADR-0015 (scope of the independent variable), ADR-0016 (the world
% is an execution model), ADR-0017 sec.1 (reserve bus), RESULTS sec.6.7-6.8.
%
% RULES FOR THIS CHAPTER:
%  * Parity is DEMONSTRATED, not asserted. Every deviation is reported WITH ITS
%    DIRECTION (EVALUATION_PROTOCOL sec.14 item 7).
%  * Minimal code naming: describe in words. Test names stay in PARITY.md.
%  * ADR-0015 rule 7 (sequential reserve) is NOT implemented. Never cite it as
%    built.

The whole thesis rests on one comparison. The agentic system and the workflow
baseline run the same evacuation, and their results are set side by side. That
comparison only means something if the two systems work under the same
conditions. If one of them sees more of the world, can say more to it, or is
scored in a different way, then the result measures the test rig and not the
architecture.

This chapter shows that the conditions are the same, and says exactly where
they are not. Parity is not claimed in prose. It is held by construction and
checked by an executable test suite. Where the two systems differ on purpose,
the difference is declared together with the direction in which it pushes the
result. The chapter comes before the results because it is what makes the
results admissible.

% =====================================================================
\section{What Parity Has to Mean Here}
\label{sec:val:definition}
% source: PARITY sec.0, sec.6

The study has one independent variable: \emph{how coordination is produced}.
On the treatment side, coordination comes from distributed LLM agents, each
holding local information and negotiating through the shared database. On the
control side, it comes from one deterministic Incident Commander (IC) that
follows written doctrine. Everything else has to be the same.

This is not a comparison of ``LLM against no LLM''. Only the two reasoning
coordination agents are replaced by the IC: the command-post agent (BIS) and
the transport agent (Hochbahn). The bus drivers and the other field agents
still run a language model in both conditions. Section~\ref{sec:val:scope}
gives the full roster.

\paragraph{Parity by identity.} A common way to build a fair control is to
write a second system that behaves like the first. That approach has a
weakness. A copy can match today and drift tomorrow, and nobody notices.
This project uses a stricter rule. The IC is built on the same engine that
all ten agents run on, and it reuses the absorbed agents' code as the very
same functions, not as copies.

Every attribute the IC defines must fall into one of three groups. Of the 49
attributes it defines, 27 are the unmodified functions of the two absorbed
agents (23 from BIS and 4 from Hochbahn). The other 22 are new, and each one is
declared with its reason. Some are merged, because both agents defined them
and the IC needed one version. One is the seam where the independent variable
itself sits. Some belong to the message link between the two agents, which no
longer exists once they are one node. Everything else is inherited unchanged
from the shared engine. The test checks identity, not behaviour. A
reimplementation that behaves the same fails the test on the day it is
written. A new method that nobody justified fails it too.

There is one small difference from the design record. The architecture
decision said the IC state should \emph{contain} the two agents' states. In
the code it \emph{flattens} them into one state. This was needed so that the
absorbed functions share one command lifecycle, one inbox and one event log.
The test checks that every field keeps its name, type and default.

\paragraph{What this costs.} An earlier version of the audit had a more vivid
result. It ran one scenario twice, once through each side, and found that the
two event logs differed in one field across 30 rows. That instrument could not
survive the merge of the two agents into the IC. The IC's vocabulary is a
superset of one agent's, so a replay through that agent drops the verbs it
does not know, and the logs would differ for reasons that are not real parity
breaks. The paired run was therefore retired. Three of the five axes below are
now checked \emph{structurally}: the test shows that the two sides use the same
code, not that they behaved alike on one run. Identity is the stronger claim,
because it holds for every scenario and not only the one that was walked. It
is also the less vivid one, and that loss is stated here.

% =====================================================================
\section{The Capability Table}
\label{sec:val:table}
% source: PARITY sec.1-5

Parity is audited on five axes. Table~\ref{tab:val:capability} summarises
them. For each axis it names what the two systems share, how the sharing is
held and checked, and where they differ on purpose.

\begin{table}[htbp]
\centering
\caption{The five parity axes. ``Identity'' means the baseline uses the same
code objects as the agentic side, not copies.}
\label{tab:val:capability}
\small
\begin{tabularx}{\textwidth}{>{\raggedright\arraybackslash}p{2.2cm}>{\raggedright\arraybackslash}X>{\raggedright\arraybackslash}p{2.3cm}>{\raggedright\arraybackslash}X}
\toprule
Axis & What is shared & How it is held & Declared difference \\
\midrule
Information horizon & Same inbox, same sensing snapshots, same simulation time. No fact reaches the IC by a route no agent has. & Identity, plus a walk through a scenario & The IC holds the command picture and the fleet picture at once (deviation~5). \\
Action vocabulary & The union of the two absorbed agents' command sets, minus the verbs of the deleted link. & Equality test against the live agents & Four verbs cannot be reached; the baseline uses 9 of its 13. \\
World model & Same sensing entry points, same time-resolution pass when an action is committed. & Identity and a source sweep & The IC plans with road distance and duration (deviation~6). \\
Scenario inputs & The same scenario files through the same loader. & Shared function, tested on every file & Message envelope details (deviations~1, 2); only the baseline files declare a world (deviation~7). \\
Output substrate & The same row format in the same \code{event_log}, read by the same scorer. & Inheritance, asserted by identity & A different row label and an empty model name (deviations~3, 4). \\
\bottomrule
\end{tabularx}
\end{table}

\paragraph{Information horizon.} Both sides receive the same messages at the
same simulation time, and both read the world through the same sensing
functions. The test walks a scenario and checks three things. The inbox holds
exactly the messages the turns delivered. Every sender is a known party. Both
snapshots come from the shared sensing code. There is one real difference.
On the agentic side, the command picture and the fleet picture live in two
agents, which must message each other to combine them. The IC holds both at
once. It sees more \emph{at once}, but it sees nothing \emph{from elsewhere}.
This is deviation~5 in Section~\ref{sec:val:deviations}.

\paragraph{Action vocabulary.} BIS has 11 coordination verbs and Hochbahn has
7. They share one, so the union is 17. Merging the two agents deletes the link
between them, and four verbs named that link only: asking for a deadline
extension, reporting a finished transport, waiting for transport, and
requesting buses. Under central command there is nothing to negotiate, because
the transport section and the commander are the same authority. The IC
therefore has 13 verbs. This difference is part of what the thesis measures,
and it is stated in advance: the agentic system can show an extension
negotiation and the baseline cannot. A metric that rewarded such a negotiation
would be measuring a link one architecture does not have.

The test asserts \emph{equality} with the union, not only that one set is
contained in the other. A verb that goes missing from the control is as much a
capability gap as an extra one. The IC holds no physical verbs. Driving stays
with the driver agents on both sides.

Availability is equal. What differs is \emph{reach}. The baseline's procedure
issues 9 of its 13 verbs. The four it never issues are two verbs that end the
mission, one that stands the escorts down for good, and a free-text message
channel that a deterministic procedure has no use for. The gap is declared, and
a test fails if it ever grows. A last check reads the agentic prompts
themselves. A worked example in a prompt could teach a model a verb the
baseline lacks. The sweep found none.

\paragraph{World model.} Each absorbed agent has exactly one place where it
reads the world, and the IC calls those two functions and nothing else. The
shared sensing back ends are seeded the same way and are not replaced during a
run. When an action is committed, both sides run the same timing and
feasibility pass, because it is the same inherited method. One difference is
new. The IC's run schedule reads real road distance and duration from the
routing service. The agentic coordinator never sees these numbers; only the
drivers do. This is deviation~6. It affects the IC's \emph{planning}, not the
commit step this axis is about. The road data is cached on disk, and a cache
miss falls back to a fixed straight-line estimate, so baseline runs stay
reproducible.

\paragraph{Scenario inputs.} Both sides load every scenario file through the
same function, and the test runs it over all 15 files. Two details of the
message envelope differ (deviations~1 and~2), and neither ever reaches a
decision. One paired scenario has a different head-count on purpose. In the
road-closure scenario, the agentic side can re-route, so the closure costs time
and distance while both shelters stay reachable, and all 100 residents are
admissible. The baseline has no routing verb. The only way to express a
closure to it is as a shelter becoming unreachable. Losing a 60-bed shelter
would leave 40 residents with no bed, so the baseline file is sized at 60.
This stays fair because completeness is scored against each scenario's own
declared population, not a raw head-count. Everything else the two runs are
scored in, such as the hazard zone, the flood arrival and the fleet, is
asserted identical. A test pins that this is the only scenario with a
non-nominal population.

\paragraph{Output substrate.} Both sides write the same seven-field row format
into the same append-only \code{event_log}, and the rows are written by the
same inherited code. The event numbering is checked to be dense across a whole
run. A gap would mean a second writer or a lost row. Two labels differ on
purpose. A baseline decision row is labelled ``decision'' instead of ``LLM
response'', and its model name is empty. Logging a model-free decision as an
LLM response would be exactly the kind of quiet mislabelling the shared log
exists to prevent. The test shows the label is harmless: relabelling every
baseline decision row leaves the facts the scorer extracts byte-identical.

% =====================================================================
\section{Scope of the Independent Variable}
\label{sec:val:scope}
% source: PARITY sec.6, ADR-0015 and its status note

A baseline run replaces every \emph{reasoning coordination} agent and nothing
else. Table~\ref{tab:val:roster} shows how the ten agents split.

\begin{table}[htbp]
\centering
\caption{Which agents the baseline replaces.}
\label{tab:val:roster}
\small
\begin{tabularx}{\textwidth}{>{\raggedright\arraybackslash}p{3.2cm}>{\raggedright\arraybackslash}X>{\raggedright\arraybackslash}p{1.6cm}>{\raggedright\arraybackslash}p{3.4cm}}
\toprule
Class & Agents & LLM? & In a baseline run \\
\midrule
Reasoning coordination & BIS (command post), Hochbahn (transport) & yes & Absorbed into the IC \\
Scripted coordination & Augustinum, shelter 1, shelter 2, disaster-management staff (ZKD) & no & Unchanged, shared by both \\
Field execution & bus driver 1, bus driver 2, police, Red Cross & yes & Unchanged, shared by both \\
\bottomrule
\end{tabularx}
\end{table}

The eight agents that are not absorbed are identical in both conditions by
construction. That is what puts them outside the comparison. The four scripted
agents are already deterministic, so moving them would change no variable.
The four field agents stay out because driving a route is the same act under
either architecture, and leaving them untouched is a stronger guarantee than
auditing them.

This has a direct consequence for how the result is read. The independent
variable is how coordination is produced: one deterministic command cycle that
holds every coordination function, against distributed LLM agents that hold
local information. It is not ``LLM against no LLM'' system-wide.

Two further points about scope must be stated. First, since the IC absorbed
both agents, the baseline plans its own evacuation. It writes a run schedule
with a passenger split, per-run deadlines and a delivery estimate that can be
checked. ``The baseline is a competent command system'' is therefore testable,
not assumed. Second, the design record for this change contains a rule that
the baseline should use the reserve bus only in sequential runs. \textbf{That
rule is not implemented}, and this thesis does not claim it as built.
Section~\ref{sec:val:reserve} describes what the code does instead. Baseline
runs from before this change behave differently and are not comparable with
the evaluated ones, so every analysis is scoped by run.

% =====================================================================
\section{The Seven Declared Deviations}
\label{sec:val:deviations}
% source: PARITY "Declared deviations - the complete list"; RESULTS sec.6.7

The audit names seven deviations and no others. Anything else that shows up as
a difference between the two logs would be a defect.
Table~\ref{tab:val:deviations} lists them with their direction.

\begin{table}[htbp]
\centering
\caption{The seven declared deviations, each with the direction in which it
can move the result.}
\label{tab:val:deviations}
\small
\begin{tabularx}{\textwidth}{c>{\raggedright\arraybackslash}X>{\raggedright\arraybackslash}p{1.6cm}>{\raggedright\arraybackslash}X>{\raggedright\arraybackslash}p{2.6cm}}
\toprule
\# & Deviation & Axis & Why it exists & Direction \\
\midrule
1 & Message ids are numbered by position, not by wall clock & Inputs & Two runs of one scenario must not differ & Inert: the id is never read \\
2 & The message envelope carries its recipient & Inputs & The offline run has no dispatcher to route it & Inert: nothing reads it \\
3 & Decision rows are labelled ``decision'', not ``LLM response'' & Output & A model-free run must not be logged as an LLM response & Inert: scored facts byte-identical \\
4 & The model name is empty & Output & It is the marker of the control condition & Inert: the scorer reads no model name \\
5 & One node holds the command and fleet pictures at once & Horizon & Central command consolidates at the command post & Favours the baseline \\
6 & The IC plans with road distance and duration & World model & A schedule in minutes needs travel times & Favours the baseline \\
7 & Only the baseline scenarios declare a world & Inputs & Removes outcomes that were written into the script & Against the baseline \\
\bottomrule
\end{tabularx}
\end{table}

The seven are of three different kinds.

\paragraph{Deviations 1 to 4 are inert.} Each one exists for a reason of
honesty or determinism, and none can reach a decision or a score. Each is
pinned by its own test. The message id, for example, is replaced by a fresh one
the moment a message enters an inbox, so the difference never reaches the
situation picture, a prompt or a decision.

\paragraph{Deviations 5 and 6 are real advantages to the control.} They follow
from the doctrine. A central commander who could not see what they command
would be a strawman, and a schedule measured in minutes needs travel times.
They were accepted for these reasons and, in the audit's own words,
``reported rather than corrected''. This cuts in the direction that costs the
thesis its preferred result, and it has to be read both ways. An agentic win
despite these deviations would have been \emph{stronger}, because distributed
reasoning would have beaten a commander with a better view. The agentic loss
that was observed is \emph{weaker} than it looks, because part of the gap may
come from the wider view and not from the coordination procedure.

\paragraph{Deviation 7 runs the other way.} Only the baseline's scenario files
declare a world, which a small simulator then plays out. The agentic files do
not yet use it. This removes an advantage the control used to have rather than
adding one. It was still open when the evaluation ran, and the next section
explains why it does not threaten the result.

% =====================================================================
\section{Deviation 7 and Why It Does Not Threaten the Result}
\label{sec:val:dev7}
% source: PARITY sec.4e; RESULTS sec.6.8 item 1; THESIS_SOURCES sec.4 item 3

The evaluation protocol set a gate before any confirmatory claim: deviation~7
had to be ``closed, or re-argued in writing''. When the campaign ran, it was
neither. The re-argument did exist, in the parity audit, but it was never
named as the discharge of this gate. It is named here.

The argument is about what the deviation removes. Before the world was
declared, the baseline's scenarios injected fixed drop-off confirmations. The
same messages arrived at the same times and named the same shelters, whatever
the doctrine had decided. The scorer never read where the residents were
actually sent. A baseline that sent every resident to a shelter that had
closed would still have scored full completeness. This was the last place
where the control's outcome was \emph{written down rather than earned}. The
agentic side never had this help, because its counterparties are live agents.
Declaring a world on the baseline side takes the help away. The deviation
therefore runs \emph{against} the baseline.

That direction settles the question for this thesis. The result is that the
\textbf{baseline won}. A deviation that disadvantages the baseline cannot
explain a baseline win. At most it made the baseline's margin look smaller
than it is.

The shape of this argument should be stated openly. It works only because of
the direction of the result. Had the agentic side won, the gate would have had
to be closed properly, by giving the agentic scenarios the same world, before
that claim could stand. The deviation is therefore declared as \emph{an
asymmetry that remained open, whose direction is opposed to the result, and
which is therefore not a threat to it}.

% =====================================================================
\section{Two Further Asymmetries Outside the Numbered Seven}
\label{sec:val:extra}
% source: PARITY sec.3d (residual) and audit item 5; RESULTS sec.6.7

Two more differences are not on the list of seven, but they belong in this
chapter and in the limitations. Both favour the baseline.

\paragraph{Where a replacement bus leg starts.} Bus kilometres are measured the
same way on both sides: a leg is charged once, when the bus starts it, from
where the bus stands. One difference remains. An agentic driver's next leg
starts from its real position, which the world reports back. The baseline has
no live position to read, so it moves the bus to the place it was heading for.
This is the cheaper of the two possible starting points, so the difference
favours the baseline. It was measured on three scenarios: 83.1\,km against
89.1\,km on the baseline road closure, 99.2 against 105.2 on shelter capacity,
and 92.8 against 95.8 on bus breakdown. That is at most 6\,km, and no level
moves anywhere in that band. Closing it would mean giving the control a
position model with no runtime to feed it, so it was declared rather than
tuned away.

\paragraph{A fallback prompt that mis-frames two message types.} The command-post
agent chooses a short instruction block, a playbook, for each kind of incoming
message. It has no playbook for a route update or a bus breakdown. Both fall
to the fallback, which is meant for messages the runtime could not classify.
Its instruction is to assume that the resource is \emph{not} coming. The
baseline handles both disruptions with written branches. This is an
agentic-side deficit, and it makes the comparison on the road-closure scenario
(S02) unfair \emph{in the baseline's favour}. The fix was due before any
agentic road-closure run was scored. That deadline passed, and the item now
stands as a disclosed limitation. One qualification limits its reach. The
agentic road-closure scenario delivers the closure as a situation update to
the transport agent, and a situation update has its own playbook. Whether the
command-post model ever met a route update in a scored run is a question for
the campaign logs. This audit does not answer it.

\paragraph{What the three non-inert asymmetries have in common.} Deviations~5
and~6 and these two further items all run in the baseline's favour. Together
they can explain part of the margin. None of them can turn the result around
or explain the agentic deficit away. The one asymmetry that runs against the
baseline, deviation~7, can only have made the margin look smaller.

% =====================================================================
\section{The Reserve-Bus Mechanism}
\label{sec:val:reserve}
% source: PARITY reserve-bus box; ADR-0017 sec.1; THESIS_SOURCES sec.4 item 2
% A NAMED MECHANISM, NOT A CONFOUND (author decision 2026-09-06).

One difference between the two systems is not a deviation at all, and it is
named here so that it is not mistaken for one.

Every scenario declares the second bus as a \emph{reserve}, with 50 seats. The
baseline's scheduler never reads that role. It gives each leg to the bus that
is free first, so the second leg goes to bus~2 at the start, and the two buses
run in parallel. The agentic transport agent sees the role and keeps bus~2
parked by default. It releases it for only four named reasons: a load larger
than bus~1 can carry, bus~1 being unavailable, the command post naming two
buses, or a deadline test that fails when the last delivery would land less
than 30 minutes before the mission deadline. Otherwise it runs bus~1 in
serial trips.

The reserve role is \emph{soft}. This was the author's decision on
2026-09-06, and it is recorded with the other decisions taken without a formal
design record. A reserve is there to be used when residents would otherwise
drown. Both sides have the same fleet and both may use bus~2. Sending it out
in parallel and keeping it parked are two \emph{decisions} over one
capability, and decisions are what the comparison measures. The mechanism is
therefore not counted among the deviations. Leaving the reserve parked while
residents are at risk is reported as a \emph{decision-quality failure} of the
agentic side. It is a named mechanism behind part of the gap, not noise.

% =====================================================================
\section{How Parity Is Verified}
\label{sec:val:verification}
% source: PARITY "On test strength", "Verification"; workflow-baseline-README l.33

Every claim in this chapter is backed by an executable test. The parity suite
has 17 test functions, which expand to 31 test items. It needs no database and
no model. The baseline package as a whole carries 160 test items. One test
closes the loop between the audit and the code: every parity test must be
cited by a row of the capability table, and every row must cite a test that
exists. A test that no row cites is a claim nobody made. A row that cites a
missing test is a stale table.

The suite also carries guards against empty passes. A check such as ``every
reused function is the original'' passes trivially if nothing is reused. Each
key assertion is therefore paired with a check that the thing it tests is
actually there, for example that at least one absorbed function is reached and
that each prompt module teaches at least five verbs.

Three weaknesses of the verification are disclosed.

\begin{itemize}
\item \textbf{No mutation testing of the rewritten suite.} The earlier suite
  was checked by deliberately breaking the code and confirming that a test
  failed. When the IC absorbed the two agents, most of those tests were
  rewritten. The same check was planned for the new tests, with a minimum set
  of nine mutations, but it was not run before the freeze of 2026-09-10.
\item \textbf{The audit was absent for a short interval.} The design record
  required the new gate to be written, and to fail, before the engine changed,
  so that it was never missing. In fact the engine changed first and the gate
  followed.
\item \textbf{One small bookkeeping error.} The list of declared new IC code
  has 23 entries, while 22 are warranted. One entry is in fact reused by
  identity, so its declaration is redundant. It is harmless, because the test
  checks reuse first.
\end{itemize}

\paragraph{What this chapter licenses.} The two systems share their world,
their inputs, their vocabulary and their output substrate, and most of that
sharing is shown by identity of code rather than by similar behaviour. Where
they differ, the difference is declared with its direction. The non-inert
differences favour the baseline, and the one that does not favour it can only
have hidden part of its margin. The results in Chapter~\ref{ch:results} can
therefore be read as a comparison of two ways of coordinating. They cannot be
read as a clean measure of the coordination procedure alone, because
deviations~5 and~6 give the control a wider view by design.
